\section{The Economic Profile of Cartel-Adjacent Losers}
\label{sec:cobidder_type}

The validation results show that the award-layer ranking
concentrates adjudication-anchored loser-side exposure. This
section asks whether that exposure target has economic content.
The comparison is intentionally demanding: among firms already
inside the frequent-loser stratum, do cobidders differ from
other frequent losers in dimensions relevant for loser-side
deployment?

The answer matters for interpretation. If cobidders look like
other frequent losers once both groups cross the same
participation threshold, the validation target would be easier
to read as mechanical exposure. If they differ in deployment
breadth, proximity to legal cartel anchors, market
specialization, and bid-level losing behavior, the target becomes
more meaningful for forensic prioritization. This profile-based
reading is consistent with empirical work showing that collusion
in auctions can appear through bidder roles, repeated
relationships, and cartel organization rather than through a
single universal price statistic. Studies of timber, school
milk, and highway or procurement cartels show role allocation and
repeated relationships in concrete bidding environments
\citep{baldwin1997bidder,porter1999ohio}; auction evidence
extends the same logic to first-price collusion
\citep{pesendorfer2000study}. Case-based and dynamic analyses
further show how internal cartel
organization shapes observed bidding behavior
\citep{asker2010leniency,clark2021collusion}. The evidence below
supports the latter interpretation.

\subsection{Comparing Cobidders to Other Frequent Losers}
\label{sec:cobidder_firm}

The firm-level comparison holds fixed the award-layer selection
margin. Both groups are always-losers, and both cross the
frequent-loser threshold. Cobidders are frequent losers observed
with direct CADE defendants; the comparison group is frequent
losers without that adjacency. The question is whether the
adjudication-anchored target marks a different loser-side
profile among firms already selected by persistent losing.

The differences line up with the triage interpretation.
Cobidders are deployed more broadly: they participate in
$\valBridgeTendCob$ tender-items, compared with
$\valBridgeTendFLnc$ for other frequent losers
($d=\valBridgeTendD$). They also face more unique winners,
$\valBridgeUniqWinCob$ versus $\valBridgeUniqWinFLnc$
($d=\valBridgeUniqWinD$). These are deployment measures, not
legal conclusions: they show that cobidders occupy a larger
loser-side footprint within the already-selected frequent-loser
pool.

They are also closer to the legal anchors used in validation.
Among cobidders, $\valBridgeDirectCob$ of loss-bids face a direct
CADE defendant, compared with $\valBridgeDirectFLnc$ for other
frequent losers ($d=\valBridgeDirectD$). Their portfolios are
more concentrated in product space as well, with item-group HHI
$\valBridgeHHICob$ against $\valBridgeHHIFLnc$
($d=\valBridgeHHID$). The combination of broader deployment and
narrower product concentration is exactly the kind of profile
that makes the target useful for forensic prioritization.

The profile is not one-dimensional. Cobidders have a lower share
of own loss-bids in repeat winner-pair clusters than other
frequent losers, $\valBridgeRepeatShareCob$ versus
$\valBridgeRepeatShareFLnc$ ($d=\valBridgeRepeatShareD$). That
wrong-signed measure is informative: the cobidder label is not a
mechanical proxy for a single winner-pair pattern. The stronger
regularities are deployment breadth, proximity to CADE anchors,
and product specialization within the frequent-loser stratum.

\subsection{Bid-Level Behavior}
\label{sec:cobidder_bid}

The LANCES export lets us ask whether the distinction also
appears inside submitted bids. This evidence uses the forensic
layer; it is deliberately kept outside the award-layer screen.
For firms with at least five usable losing bids in convite or
preg\~ao, we compute the gap between each losing bid and the
winning bid in the same tender-item,
\[
  \frac{\text{bid}-\text{winning bid}}{\text{winning bid}},
\]
winsorized at $[-0.99,10]$ before firm-level aggregation. The
usable sample contains $\valTBBLnCob$ cobidders and
$\valTBBLnFLnc$ frequent-loser non-cobidders.

Two patterns matter. First, cobidders bid closer to winners:
their median bid gap is $\valTBBLMedGapCob$, compared with
$\valTBBLMedGapFLnc$ for other frequent losers
($d=\valTBBLMedGapD$). In a cover-bidding environment, losing
bids need not be visibly extreme; plausible losing bids can
preserve the appearance of rivalry. Second, cobidders show
higher within-firm dispersion in bid gaps,
$\valTBBLSDGapCob$ versus $\valTBBLSDGapFLnc$
($d=\valTBBLSDGapD$), consistent with firms appearing in
different losing roles across tenders.

The bid-level distinction is not just participation intensity in
another form. In a logit of cobidder status within the
frequent-loser stratum, controlling for
$\log(1+\text{tenders\_count})$, both signs remain: the
median-gap coefficient is $\valTBBLLogitMedCoef$
($z=\valTBBLLogitMedZ$), and the dispersion coefficient is
$\valTBBLLogitSDCoef$ ($z=\valTBBLLogitSDZ$). The forensic layer
therefore gives independent economic content to the
adjudication-anchored loser-side target.

\begin{table}[htbp]
\centering
\caption{Cobidder profile within the frequent-loser stratum}
\label{tab:cobidder_signature_submission}
\footnotesize
\begin{adjustbox}{max width=\textwidth,center}
\begin{tabular}{p{4.4cm}ccc}
\toprule
Dimension & Cobidders & FL non-cobidders & Difference \\
\midrule
Total participations & $\valBridgeTendCob$ & $\valBridgeTendFLnc$ & $d=\valBridgeTendD$ \\
Unique winners faced & $\valBridgeUniqWinCob$ & $\valBridgeUniqWinFLnc$ & $d=\valBridgeUniqWinD$ \\
Loss-bids facing direct CADE defendant & $\valBridgeDirectCob$ & $\valBridgeDirectFLnc$ & $d=\valBridgeDirectD$ \\
Portfolio item-group HHI & $\valBridgeHHICob$ & $\valBridgeHHIFLnc$ & $d=\valBridgeHHID$ \\
Repeat winner-pair share & $\valBridgeRepeatShareCob$ & $\valBridgeRepeatShareFLnc$ & $d=\valBridgeRepeatShareD$ \\
\midrule
Median bid gap to winner & $\valTBBLMedGapCob$ & $\valTBBLMedGapFLnc$ & $d=\valTBBLMedGapD$ \\
Within-firm SD of bid gap & $\valTBBLSDGapCob$ & $\valTBBLSDGapFLnc$ & $d=\valTBBLSDGapD$ \\
\bottomrule
\end{tabular}
\end{adjustbox}
\end{table}

\subsection{Interpretation}
\label{sec:cobidder_interpretation}

The profile gives economic content to the validation target.
Within the frequent-loser stratum, cobidders are more broadly
deployed, closer to legal cartel anchors, more concentrated in
product space, closer to winners in bid gaps, and more dispersed
across their own losing bids. These are precisely the dimensions
on which a forensic queue should differ if the award-layer rank
is capturing more than generic losing volume.

This evidence also clarifies the division of labor between data
layers. The award layer identifies a priority set using
information available before bid recovery. The LANCES bid layer
then shows that the priority set has economically meaningful
loser-side behavior. That is the paper's staged logic in
miniature: cheap records order attention, and richer records
evaluate what the screen has prioritized.

The next section turns from interpretation to implementation. It
asks whether the award-layer rank can be placed before
bid-distribution forensics in an operational sequence that saves
bid-microdata recovery while preserving adjudication-anchored
recall.
